\documentclass[11pt]{article}
\usepackage[margin=1in]{geometry}
\usepackage{amsmath,amssymb,amsthm}
\usepackage{graphicx}
\usepackage[hidelinks]{hyperref}
\setlength{\emergencystretch}{2em}

\newtheorem{theorem}{Theorem}
\newtheorem{lemma}{Lemma}
\newtheorem{definition}{Definition}

\newcommand{\C}{\mathbb C}
\newcommand{\Kcal}{\mathcal K}
\newcommand{\Prim}{\operatorname{Prim}}
\newcommand{\supp}{\operatorname{supp}}
\newcommand{\Ball}{\operatorname{Ball}}

\title{Unitizations of Separable-Fibre Continuous-Trace Algebras Have BCP}
\author{Run \texttt{fa\_banach\_001}, lane 19}
\date{June 17, 2026}

\begin{document}
\maketitle

\section*{Source Problem}

The source paper is Jinghao Huang, Karimbergen Kudaybergenov, and Rui Liu,
\emph{The ball-covering property of von Neumann algebras and noncommutative
symmetric spaces}, arXiv:2311.14257.  After proving that a commutative unital
\(C^*\)-algebra \(C(K)\) has BCP/UBCP exactly when its pure state space, i.e.
\(K\), has a countable \(\pi\)-basis, the authors ask for a noncommutative
version:

\begin{center}
\includegraphics[width=0.92\textwidth]{figures/open_problem_crop.png}
\end{center}

\section*{Statement}

The previous continuous-trace packet proves the expected criterion for
continuous-trace algebras whose elementary fibres are separable.  The present
packet proves that the positive theorem survives one scalar multiplier
direction.

\begin{theorem}[unitized continuous-trace subcase]
\label{thm:main}
Let \(I\) be a continuous-trace \(C^*\)-algebra whose elementary fibres are
\(\Kcal(H_x)\) with \(H_x\) separable and nonzero.  Let \(A=\widetilde I\) be
the minimal unitization of \(I\) when \(I\) is nonunital; if \(I\) is unital,
put \(A=I\).  Then
\[
        A\text{ has BCP}
        \quad\Longleftrightarrow\quad
        A\text{ has UBCP}
        \quad\Longleftrightarrow\quad
        \Prim(I)\text{ has a countable }\pi\text{-basis}.
\]
\end{theorem}

\section*{Idea of the Proof}

The old compact-bump proof works because local compact-operator sections can
be multiplied by scalar bumps.  In a unitization, an element has local form
\[
        a(x)=\lambda I+i(x),
\]
and \(\lambda I\) cannot be localized when the fibre is infinite-dimensional.
This is the new obstacle.

The trick is to use the scalar direction as margin rather than trying to
localize it.  If \(|\lambda|\) is tiny, the compact part \(i(x)\) nearly norms
the element and the old compact proof still works.  If \(|\lambda|\) is
bounded below, choose a finite-rank approximation \(T\) to the local compact
part and add a large rank-one compact ``ballast'' \(M\omega Q\) in a direction
orthogonal to \(T\).  The center is
\[
        c=\mu 1+h(T+M\omega Q),
\]
where \(\mu\) approximates \(\lambda\), \(\omega\) has the same phase, and
\(h\) is a local bump.  On the spike direction, \(\|c\|\) is about
\(M+|\lambda|\), while the distance from \(a\) to \(c\) is only about \(M\).
That fixed scalar margin is exactly what the naive localization argument was
missing.

\section*{Proof}

We use two standard facts already used in the companion packets.

\begin{lemma}[primitive \(\pi\)-base obstruction]
\label{lem:primitive}
If a \(C^*\)-algebra \(B\) has BCP, then \(\Prim(B)\) has a countable
\(\pi\)-basis.
\end{lemma}

\begin{lemma}[charted countable bump families]
\label{lem:charts}
Assume that \(X=\Prim(I)\) has a countable \(\pi\)-basis.  There are countably
many continuous-trace trivialization charts
\[
        \Phi_n:I_{O_n}\longrightarrow C_0(O_n,\Kcal(H_n)),
\]
where \(H_n\) is separable and nonzero, such that \((O_n)\) is a
\(\pi\)-basis of \(X\).  For each \(n\), there is a countable family
\(\mathcal B_n\subset C_c(O_n)\) with \(0\le h\le1\), \(\|h\|_\infty=1\),
and such that every nonempty open \(W\subset O_n\) contains the support of
some \(h\in\mathcal B_n\).
\end{lemma}

\begin{proof}
This is the chart-and-bump selection from the continuous-trace packet.  Local
triviality of continuous-trace algebras gives charted neighbourhoods.  Refining
the given countable \(\pi\)-basis by these charted neighbourhoods gives a
countable chart \(\pi\)-basis.  Since each \(O_n\) is locally compact
Hausdorff, compactly supported scalar bumps with cofinal supports are obtained
from local regularity.
\end{proof}

\begin{proof}[Proof of Theorem~\ref{thm:main}]
If \(I\) is unital, then \(A=I\), and the result is exactly the previous
continuous-trace separable-fibre theorem.  We therefore assume that \(I\) is
nonunital and write elements of \(A\) uniquely as
\[
        a=\lambda 1+i,\qquad \lambda\in\C,\quad i\in I.
\]

The implication UBCP \(\Rightarrow\) BCP is immediate.  If \(A\) has BCP,
Lemma~\ref{lem:primitive} gives a countable \(\pi\)-basis for \(\Prim(A)\).
The open subset \(\Prim(I)\subset\Prim(A)\) then also has a countable
\(\pi\)-basis.  It remains to prove UBCP assuming that \(\Prim(I)\) has a
countable \(\pi\)-basis.

Fix the charts and bump families from Lemma~\ref{lem:charts}.  We shall build
countably many balls covering \(S_A\) and uniformly avoiding a neighbourhood of
the origin.

First include the old compact-bump centers.  For each \(n\), choose a
countable \(1/64\)-net \(D_n\subset S_{\Kcal(H_n)}\).  For \(h\in\mathcal B_n\)
and \(T\in D_n\), include
\[
        c_0(n,h,T)=\Phi_n^{-1}(4hT)\in I\subset A
\]
with ball radius \(7/2\).  These centers have norm \(4\).

Next include scalar-ballast centers for the infinite-dimensional charts.  Put
\[
        \eta=\frac1{16},\qquad \varepsilon=\frac1{256},\qquad M=4.
\]
For each infinite-dimensional \(H_n\), fix an orthonormal basis
\((e_m)_{m\ge1}\).  Let \(\mathcal F_n\) be the countable set of finite-rank
operators whose matrices, relative to this basis, have rational complex
entries and finite support.  This set is norm dense in \(\Kcal(H_n)\).  Let
\(\mathcal Q_n=\{Q_m:m\ge1\}\), where \(Q_m\) is the rank-one projection onto
\(\C e_m\).  Let \(\Lambda\subset\C\) be a countable dense set and let
\(\Omega\subset\mathbb T\) be a countable dense set.  For every
\[
        n,\ h\in\mathcal B_n,\ T\in\mathcal F_n,\ Q\in\mathcal Q_n,\ 
        \mu\in\Lambda,\ \omega\in\Omega
\]
include the center
\[
        c_1(n,h,T,Q,\mu,\omega)=
        \mu1+\Phi_n^{-1}\bigl(h(T+M\omega Q)\bigr).
\]
We retain in the covering family only those of these countably many centers
which satisfy
\[
        \|c_1\|\ge M+\eta-2\varepsilon,
\]
and for them use the fixed radius \(M+3\varepsilon\).  Since
\[
        M+\eta-2\varepsilon-(M+3\varepsilon)=\eta-5\varepsilon>0,
\]
these balls are uniformly off the origin.

We now prove that the two families cover \(S_A\).  Let \(a=\lambda1+i\in S_A\).
For \(P\in O_n\), write
\[
        i_n(P)=\Phi_n(i|_{I_{O_n}})(P)\in\Kcal(H_n).
\]

Suppose first that either \(|\lambda|<\eta\), or the selected chart below has
finite-dimensional fibre.  The set
\[
        W=\{P\in\Prim(I):\|\lambda I+i(P)\|>3/4\}
\]
is nonempty and open; if the norm of \(a\) is attained only at the quotient
character, then \(|\lambda|=1\), so this case cannot occur.  Choose a chart
\(O_n\subset W\) and \(P_0\in O_n\).

If \(H_n\) is finite dimensional, set
\[
        b(P)=\lambda I_{H_n}+i_n(P)\in \Kcal(H_n).
\]
If \(H_n\) is infinite dimensional, set \(b(P)=i_n(P)\).  In the infinite
case,
\[
        \|b(P_0)\|\ge \|\lambda I+i_n(P_0)\|-|\lambda|
        >\frac34-\eta>\frac{11}{16}.
\]
In the finite-dimensional case, \(\|b(P_0)\|>3/4\).  Thus in both cases
\(\alpha=\|b(P_0)\|>11/16\).  Choose \(T\in D_n\) with
\[
        \left\|T-\frac{b(P_0)}{\alpha}\right\|<\frac1{64}.
\]
For \(0\le s\le1\), convexity gives
\[
        \|b(P_0)-4sT\|
        \le
        \max\left\{\alpha,\ 4-\alpha(1-1/64)\right\}
        <\frac72-\eta.
\]
By norm-continuity of the compact section \(b\), there is a nonempty open
neighbourhood \(N\subset O_n\) of \(P_0\) such that
\[
        \sup_{0\le s\le1}\|b(P)-4sT\|<\frac72-\frac{\eta}{2}
        \qquad(P\in N).
\]
Choose \(h\in\mathcal B_n\) with \(\supp h\subset N\), and put
\(c=c_0(n,h,T)\).  If \(H_n\) is finite dimensional, then on \(\supp h\)
\[
        \|a(P)-c(P)\|<7/2.
\]
If \(H_n\) is infinite dimensional, then on \(\supp h\)
\[
        \|a(P)-c(P)\|
        \le |\lambda|+\|i_n(P)-4h(P)T\|
        <\eta+\frac72-\frac{\eta}{2}<\frac72.
\]
Outside \(\supp h\), the center vanishes on primitive quotients over
\(\Prim(I)\), so the distance is at most \(\|a\|=1\); at the quotient
character the distance is \(|\lambda|\le1\).  Hence \(a\) is in one of the
compact-bump balls.

It remains to handle the case in which \(|\lambda|\ge\eta\) and the chosen
chart has infinite-dimensional fibre.  Choose a nonempty open set
\[
        W=\{P\in\Prim(I):\|\lambda I+i(P)\|>3/4\};
\]
this is again nonempty because \(|\lambda|\ge\eta\) and \(\|a\|=1\).  Choose
an infinite-dimensional chart \(O_n\subset W\) if possible.  If every chart
inside \(W\) has finite-dimensional fibre, the previous paragraph applies, so
we may assume \(H_n\) is infinite dimensional.  Pick \(P_0\in O_n\) and put
\(k_0=i_n(P_0)\).

Choose \(\mu\in\Lambda\) with \(|\mu-\lambda|<\varepsilon\).  Then
\(|\mu|\ge\eta-\varepsilon\).  Choose \(\omega\in\Omega\) such that
\[
        |M\omega+\mu|\ge M+|\mu|-\varepsilon.
\]
Choose \(T\in\mathcal F_n\) with \(\|T-k_0\|<\varepsilon\).  Since \(T\) has
finite support relative to \((e_m)\), choose \(Q=Q_m\in\mathcal Q_n\) with
\[
        TQ=QT=0.
\]
By norm-continuity of \(i_n\), choose a nonempty open neighbourhood
\(N\subset O_n\) of \(P_0\) such that
\[
        \|i_n(P)-k_0\|<\varepsilon\qquad(P\in N).
\]
Choose \(h\in\mathcal B_n\) with \(\supp h\subset N\), and put
\[
        c=c_1(n,h,T,Q,\mu,\omega).
\]

Since \(\|h\|_\infty=1\), there is \(P_1\in\supp h\) with \(h(P_1)=1\).  At
that point, the local value of \(c\) is
\[
        \mu I+T+M\omega Q.
\]
Because \(TQ=QT=0\), the one-dimensional range of \(Q\) reduces this operator
and carries the scalar \(M\omega+\mu\).  Therefore
\[
        \|c\|\ge |M\omega+\mu|
        \ge M+|\mu|-\varepsilon
        \ge M+\eta-2\varepsilon.
\]

We estimate \(\|a-c\|\).  First let \(P\in\supp h\) and set \(s=h(P)\).  Write
\[
        e(P)=i_n(P)-T.
\]
Then \(\|e(P)\|<2\varepsilon\), and the local difference equals
\[
        (\lambda-\mu)I+(1-s)T-sM\omega Q+e(P).
\]
The part without \(e(P)\) is block diagonal for
\[
        H_n=QH_n\oplus(1-Q)H_n,
\]
because \(TQ=QT=0\).  On \(QH_n\) its norm is at most \(sM+\varepsilon\), and
on \((1-Q)H_n\) its norm is at most
\[
        \varepsilon+(1-s)\|T\|
        \le \varepsilon+\|k_0\|+\varepsilon
        \le 2+2\varepsilon.
\]
Thus, since \(M=4\),
\[
        \|a(P)-c(P)\|\le M+\varepsilon+2\varepsilon=M+3\varepsilon.
\]

If \(P\in\Prim(I)\setminus\supp h\), then the compact part of \(c\) vanishes
at \(P\), and
\[
        \|a(P)-c(P)\|
        =
        \|(\lambda-\mu)I+i(P)\|
        \le \varepsilon+\|i(P)\|
        \le \varepsilon+\|\lambda I+i(P)\|+|\lambda|
        \le 2+\varepsilon
        <M+3\varepsilon.
\]
At the quotient character of the unitization, the distance is
\(|\lambda-\mu|<\varepsilon\).  Taking the supremum over primitive quotients of
\(A\) gives
\[
        \|a-c\|<M+3\varepsilon.
\]
Thus \(a\) is covered by one of the scalar-ballast balls.

We have covered \(S_A\) by countably many balls.  The compact-bump balls have
centers of norm \(4\) and radius \(7/2\), hence avoid \(B(0,1/2)\).  The
scalar-ballast balls have centers at distance at least
\(\eta-5\varepsilon>0\) beyond their radius.  Therefore the whole cover is
uniformly off the origin, and \(A\) has UBCP.
\end{proof}

\section*{Scope and Limitations}

This is a genuine extension of the continuous-trace separable-fibre theorem,
but it is still only a partial result for the source question.  It handles one
scalar multiplier direction over a positive continuous-trace ideal.  It does
not yet prove permanence for arbitrary finite-dimensional quotients
\[
        0\to I\to A\to F\to0
\]
or for general type-I/Fell algebras, where matrix-valued multiplier directions
and local compression continuity have to be controlled.

The proof suggests the next attack: replace the scalar spike by a
finite-dimensional matrix-block spike.  If that works, finite-dimensional
extensions of the positive continuous-trace class should follow.

\section*{Novelty and Search Bounds}

The run indexes were searched for arXiv:2311.14257, unitization,
finite-dimensional extension, continuous trace, compact operators, scalar
multiplier, BCP, and UBCP.  Web searches on June 17, 2026 for BCP permanence,
finite-codimensional BCP, unitizations, and Fell/type-I local continuous-trace
structure found the source paper, the vector-valued \(C_0(K,X)\) stability
paper arXiv:2111.04921, Calkin/operator-space BCP papers arXiv:2412.07137 and
arXiv:2507.02261, and Fell-algebra structure literature arXiv:1207.2540, but
no exact finite-dimensional extension or unitized continuous-trace BCP theorem.

\section*{Human Review Focus}

The main proof points to check are:

\begin{itemize}
\item the scalar-ballast estimate in the \(|\lambda|\ge1/16\) case;
\item the choice of a finite-support \(T\) and a rank-one \(Q\) with
      \(TQ=QT=0\);
\item the use of the quotient character of the unitization in the norm
      estimate;
\item the finite-dimensional fibre branch, where the scalar direction is
      compact and the old continuous-trace proof applies.
\end{itemize}

If these checks pass, this packet should be treated as a new positive partial
result and as the template for the next finite-dimensional quotient push.

\begin{thebibliography}{9}

\bibitem{HKL}
Jinghao Huang, Karimbergen Kudaybergenov, and Rui Liu,
\emph{The ball-covering property of von Neumann algebras and noncommutative
symmetric spaces},
arXiv:2311.14257, 2023.

\bibitem{lane19ct}
Run \texttt{fa\_banach\_001}, lane 19,
\emph{BCP for Continuous-Trace \(C^*\)-Algebras with Separable Fibres},
solution packet
\texttt{2311.14257\_continuous\_trace\_separable\_fiber\_bcp}, 2026.

\bibitem{lane19primitive}
Run \texttt{fa\_banach\_001}, lane 19,
\emph{A Primitive-Ideal Pi-Base Obstruction for \(C^*\)-Algebras with BCP},
solution packet
\texttt{2311.14257\_primitive\_pi\_base\_obstruction}, 2026.

\bibitem{Dixmier}
J. Dixmier,
\emph{\(C^*\)-Algebras},
North-Holland, 1977.

\bibitem{RW}
I. Raeburn and D. P. Williams,
\emph{Morita Equivalence and Continuous-Trace \(C^*\)-Algebras},
American Mathematical Society, 1998.

\end{thebibliography}

\end{document}
